% Source: https://artofproblemsolving.com/community/c3h1390147_slope_derivation
% Stable identity: aops-topic-1390147-q01
Let $\ell$ be a non-vertical line, and let $P_1(x_1,y_1)$ and $P_2(x_2,y_2)$ be two distinct points on it. A student notices that the right triangle with legs $x_2-x_1$ and $y_2-y_1$ has hypotenuse $\sqrt{(x_2-x_1)^2+(y_2-y_1)^2}$, and asks why the slope of $\ell$ is defined as the \emph{ratio} $\dfrac{y_2-y_1}{x_2-x_1}$ instead.
\begin{enumerate}
\item[(a)] Prove that the ratio $\dfrac{y_2-y_1}{x_2-x_1}$ does not depend on which two distinct points of $\ell$ are used, so it measures the steepness of $\ell$, while the hypotenuse is the distance $P_1P_2$ and depends on the points.
\item[(b)] For the line through $(1,2)$ and $(4,6)$, find the slope and the distance between these two points.
\end{enumerate}
